\section{A quantum particle in electromagnetic fields (11/13)}

\subsection{Probability conservation}
For
\begin{equation}
\imth\hbar\partial_t\psi
=-\frac{\hbar^2}{2m}\nabla^2\psi+V\psi,
\end{equation}
the probability density and current are
\begin{equation}
\rho=|\psi|^2,
\qquad
\bm j=\frac{\imth\hbar}{2m}
\left(\psi\bm\nabla\psi^*-\psi^*\bm\nabla\psi\right)
=\frac{\hbar}{m}\im(\psi^*\bm\nabla\psi).
\end{equation}
They satisfy
\begin{equation}
\boxed{\partial_t\rho+\bm\nabla\cdot\bm j=0.}
\end{equation}
Hence
\begin{equation}
\frac{\dif}{\dif t}\int_R\rho\dif^3r
=-\oint_{\partial R}\bm j\cdot\dif\bm S.
\end{equation}
For charge $e$ particles, the electric charge density and electric current are given by
\begin{equation}
\rho_e=e\rho,
\qquad
\bm j_e=e\bm j.
\end{equation}

\subsection{Gauge fields and minimal coupling}
The electromagnetic fields are
\begin{equation}
\bm B=\bm\nabla\times\bm A,
\qquad
\bm E=-\bm\nabla\phi-\partial_t\bm A.
\end{equation}
Gauge transformations are
\begin{equation}
\phi\longmapsto\phi+\partial_t f,
\qquad
\bm A\longmapsto\bm A-\bm\nabla f.
\end{equation}
The Hamiltonian is
\begin{equation}
\boxed{
\hat H=\frac{[\hat{\bm p}-e\bm A(\hat{\bm r},t)]^2}{2m}
+e\phi(\hat{\bm r},t).}
\end{equation}
Gauge invariance requires
\begin{equation}
\psi\longmapsto e^{-\imth ef/\hbar}\psi,
\qquad
\psi^*\longmapsto e^{\imth ef/\hbar}\psi^*.
\end{equation}
The canonical momentum is not gauge invariant:
\begin{equation}
\hat{\bm p}=-\imth\hbar\bm\nabla
\longmapsto\hat{\bm p}-e\bm\nabla f.
\end{equation}
Instead the kinetic momentum (also known as mechanical momentum)
\begin{equation}
\boxed{\hat{\bm\pi}=\hat{\bm p}-e\bm A}
\end{equation}
is gauge invariant, and different components of the kinetic momenta generally don't commute:
\begin{equation}
\boxed{\comm{\hat\pi_a}{\hat\pi_b}
=\imth\hbar e\epsilon_{abc}B_c.}
\end{equation}

\subsection{Ehrenfest equations}
In the Heisenberg picture,
\begin{equation}
\imth\hbar\frac{\dif\hat O}{\dif t}=\comm{\hat O}{\hat H}.
\end{equation}
For the position operator,
\begin{equation}
\frac{\dif\langle\hat{\bm r}\rangle}{\dif t}
=\frac{\langle\hat{\bm\pi}\rangle}{m}.
\end{equation}
For the kinetic momentum,
\begin{equation}
\frac{\dif\langle\hat{\bm\pi}\rangle}{\dif t}
=e\langle\bm E\rangle
+\frac{e}{2m}
\left\langle\hat{\bm\pi}\times\bm B
-\bm B\times\hat{\bm\pi}\right\rangle.
\end{equation}
For a uniform and static electromagnetic field, we therefore have the following equation of motion:
\begin{equation}
\boxed{
 m\frac{\dif^2\langle\hat{\bm r}\rangle}{\dif t^2}
=e\bm E+e\frac{\dif\langle\hat{\bm r}\rangle}{\dif t}
\times\bm B.}
\end{equation}

\subsection{Solving Schrodinger equation in a uniform magnetic field}

\subsubsection{Landau gauge}
Let
\begin{equation}
\bm B=(0,0,B),
\qquad
\bm A=(0,Bx,0).
\end{equation}
Neglecting free motion along $z$,
\begin{equation}
\hat H=\frac{\hat p_x^2+(\hat p_y-eB\hat x)^2}{2m}.
\end{equation}
Since
\begin{equation}
\comm{\hat p_y}{\hat H}=0,
\end{equation}
write
\begin{equation}
\hat p_y=\hbar k_y.
\end{equation}
Then
\begin{equation}
\hat H=\frac{\hat p_x^2}{2m}
+\frac{(eB)^2}{2m}
\left(\hat x-\frac{\hbar k_y}{eB}\right)^2.
\end{equation}
Define
\begin{equation}
\omega_c=\frac{|eB|}{m},
\qquad
\ell_B=\sqrt{\frac{\hbar}{|eB|}}.
\end{equation}
The Landau levels are
\begin{equation}
\boxed{E_n=\hbar\omega_c\left(n+\frac12\right),
\qquad n=0,1,2,\ldots.}
\end{equation}
The eigenstates are
\begin{equation}
\psi_{n,k_y}(x,y)
\propto e^{\imth k_yy}
\varphi_n\left(x-\frac{\hbar k_y}{eB}\right).
\end{equation}
For $n=0$,
\begin{equation}
\psi_{0,k_y}(x,y)
\propto e^{\imth k_yy}
\exp\left[-\frac{(x-\hbar k_y/eB)^2}{2\ell_B^2}\right].
\end{equation}

\subsubsection{Landau-level degeneracy}
With periodic boundary conditions along $y$,
\begin{equation}
k_y=\frac{2\pi j}{L_y}.
\end{equation}
The guiding-center spacing is
\begin{equation}
\Delta X=\frac{\hbar\Delta k_y}{|eB|}
=\frac{2\pi\hbar}{|eB|L_y}.
\end{equation}
The degeneracy in area $L_xL_y$ is
\begin{equation}
\boxed{D=\frac{|eB|L_xL_y}{2\pi\hbar}
=\frac{\Phi}{\Phi_0},}
\end{equation}
where
\begin{equation}
\Phi_0=\frac{2\pi\hbar}{|e|}
\end{equation}
is the flux quantum.

% \subsubsection{Guiding-center coordinates}
% The guiding-center coordinates are defined as
% \begin{equation}
% \hat{\bm R}=\hat{\bm r}
% -\frac{\hat{\bm z}\times\hat{\bm\pi}}{eB}.
% \end{equation}
% They are conserved observables
% \begin{equation}
% \comm{\hat{\bm R}}{\hat H}=0,
% \end{equation}
% and within a Landau level (see HW10)
% \begin{equation}
% \hat P_n\hat{\bm r}\hat P_n=\hat{\bm R}.
% \end{equation}
% The components do not commute:
% \begin{equation}
% \comm{\hat R_x}{\hat R_y}
% =-\frac{\imth\hbar}{eB}.
% \end{equation}

\subsubsection{Filling factor}
The filling factor is defined as the ratio of total electron number and Landau level degeneracy:
\begin{equation}
\nu=\frac{N_e}{D}
=\frac{\rho_e\Phi_0}{B}.
\end{equation}
Physical observables such as conductivity and magntization oscillate with the magnetic field, known as the quantum oscillations. The quantum oscillations are periodic in $1/B$:
\begin{equation}
\boxed{\Delta\left(\frac1B\right)
=\frac{1}{\rho_e\Phi_0}.}
\end{equation}


